\aside{
    Don't you think it's time for an epigraph?
}

\wideepigraph

\epigraph{
    Nobody stays in this valley except by a full, conscious choice based on a full, conscious knowledge of every fact involved in his decision. Nobody stays here by faking reality in any manner whatever.
}{Ayn Rand (Atlas Shrugged)}

\aside{
    Or, perhaps you might prefer something from Karl Popper?
    Very well:
}

\epigraph{
    On the level of scientific discovery two new aspects emerge. The most important one is that scientific theories can be formulated linguistically, and that they can even be published. Thus they become objects outside of ourselves: objects open to investigation. As a consequence, they are now open to \emph{criticism}. Thus we can get rid of a badly fitting theory before the adoption of the theory makes us unfit to survive. \emph{By criticizing our theories we can let our theories die in our stead.} This is of course immensely important.
}{Karl Popper (The Myth of the Framework)}

\aside{
    Let's proceed.
}

\subsubsubsection{Simplex Security}

\label{sec:simplex-security}

Simplexes are secure if attacking a single simplex-chain is as difficult as attacking the whole simplex (the network).

Say the attacker controls $q$ proportion of the network's work-generation capacity (e.g., hash-rate for PoW chains), and honest nodes control $p$ proportion, such that $p+q=1$.
The attacker should only be able to perform a doublespend if $q > p$.

Successfully attacking a single blockchain requires an attacker to publish an alternate chain-segment that the fork rule considers heavier than the corresponding public chain-segment.
Those two segments will have a common ancestor, so the weight of the attacker's segment must be greater than the weight of the public chain-segment (from that common ancestor on).

Simplex-chains evaluate block-weight as the sum of work done on that block, plus work done on reflecting blocks.
So effecting a doublespend on one simplex-chain requires generating a chain-segment with more total block-weight (including reflections) than the public chain-segment.

It is not viable for the attacker to mine the attacking chain-segment in public (see \autoref{sec:dos-and-dags}), so they must mine it in private.
Blocks mined in private will not gain reflections from honest miners on other simplex-chains, so any reflections contributing to the doublespend must be created by the attacker.
The attacker's reflecting blocks (which reflect the attacking chain-segment) on other simplex-chains can be public, but the attacker's task is simpler and not detectable if the attacker mines reflecting blocks in private.\footnote{
    Additionally, if reflecting chains maintain projections as headers-only chains (i.e., construct a main chain via SPV rules), then it should be much harder for an attacker to succeed if reflections are mined in public.
}
There is no viable way for the attacker to prevent honest miners from reflecting the honest chain-segment, including new blocks that extend it, nor to prevent the honest chain-segment from reflecting blocks from other simplex-chains (again, see \autoref{sec:dos-and-dags}).

The honest chain-segment, on the targeted simplex-chain, will not be reflected by the attacker (since that would be self-defeating).
Similarly, the attacker's chain-segment will not be reflected by the honest network, since it's being mined in private.

\todo{need to add some stuff in here -- a bit more explanation and linking, particularly consider editing 'it is not viable' paragraph}

Let $r$ be the network-wide rate-of-work (hash-rate).
Over some attack duration $d$, we expect the attacker's chain-segment to weigh $qrd$, and the honest chain-segment to weigh $prd$.
Thus, for the attacker's chain-segment to win, it must be that $qrd > prd \implies q > p$.
$\blacksquare$

\subsubsubsection{Confirmation Equivalence Conjecture}

\defineTerm{Confirmation Equivalence Conjecture (CEC)}{
    The conjecture that, when using PoR and appropriately converting work,
    confirmations of reflecting chains
    can be treated as \emph{equivalent} to local confirmations of the same weight
}

The Confirmation Equivalence Conjecture is why PoR works.
If it were false, then PoR could not work, and neither would simplexes.
Additionally, this is the root of UT's $O(c)$ confirmation rate -- the conjecture implies that confirmation rate is proportional to the number of mutually reflecting chains.

Can we test it? If so, how?

Consider a traditional PoW blockchain (e.g., Bitcoin, Eth1, etc).
It's well known that the risk of a doublespend against a particular block is related to the number of \emph{confirmations} that block has.\footnote{
    See \href{https://web.archive.org/web/20220209100515/https://cloudflare-ipfs.com/ipfs/QmNUWmY94QUievK8ptoxsPyAQUsKvx1cjRyCgPcfmysAVv}{Analysis of hashrate-based double-spending}.
}
However, this is not a linear relationship; rather, it's similar to \emph{exponential decay}.
After the first few confirmations, each additional confirmation reduces the risk of a doublespend by approximately the same factor.\footnote{
    The exact factor depends on the attacker's hash-rate as a proportion of the network's (i.e., $q$).
}
So \emph{security takes \textbf{time}}, because confirmations take time.

Now, consider a simplex: the equivalence conjecture says that it doesn't matter whether confirmations come from the \emph{local} chain, or if they come from a \emph{reflecting} chain.
So, a simplex-chain in a 2-chain simplex (wherein each simplex-chain has $\mathbb{C}^\prime = 2 \times B_f$) \emph{should}, after $C$ confirmations, be as secure as a standalone blockchain with $2C$ confirmations.
Or, more generally, a simplex-chain in an $N$-chain simplex will, after $C$ \emph{local} confirmations, be as secure as a standalone blockchain after $NC$ confirmations.

This is something we can test!
(See \autoref{sec:por-simulation-results} for that.)

As blockchain architects, if the CEC is true, then UT allows us to maintain security by trading \emph{waiting time} for \emph{more simplex-chains!}

\aside{
    If you were looking for an essence of how UT breaks the core conflict of Buterin's Trilemma\dots
}

\subsubsubsection{Generalizing Doublespends}

With a traditional blockchain, the network has no "memory" of work (hashes) that have been done but are not accounted for in the main chain.
Specifically, the network gains no knowledge of how much work has been done \emph{between} block.
This is almost self-evident: each block is the \emph{sole} way that work can be \emph{added} to the chain.
Naturally, there is no smaller unit of contributed work than an individual block -- so there's no "memory".\footnote{
    We're not concerned with things like \emph{weak blocks} or super-block hybrid designs; though it's possible that doublespend methodology for those designs has some overlap with this section.
}

When an attacker publishes a heavier chain-segment which causes a reorganization, it is \emph{immediately} in the interests of miners to begin building on the attacker's best block.
A heavier chain-segment is \emph{decisively better} in all cases.
That may seem so obvious that it isn't worth mentioning, but it's actually a \emph{specific case} that only works for traditional blockchains.
If there were some "memory", it might be worth miners \emph{continuing} mining their \emph{existing} block, \emph{instead} of reorganizing and building on the attacker's best block.

What kind of "memory" could there be?
Are \emph{Proofs of Reflection} a "memory"?

As it turns out: yes.

Let the total chain-weight of the attacker's best block be called $W_A$, and the total chain-weight of the honest network's best block be called $W_H$.

Consider a miner of a simplex-chain (chain $L$) who is working on a \emph{draft block} during an active attack (though, of course, the miner does not know that).
Specifically, let's consider \emph{just before} the attacker publishes their chain-segment.
In the memory pool\footnote{
  The \emph{memory pool} is where unconfirmed transactions (and other things) are stored before they are included in a block.
  Each node has their own memory pool -- this must be the case because securely synchronizing memory pools between untrusted parties is as difficult as running a whole blockchain.
  A miner typically decides a block's contents by choosing the combination of transactions from their memory pool that results in the largest reward (i.e., transactions with the largest fees).
} of that miner, there will be PoRs, and those PoRs are \emph{only} valid for the \emph{honest} chain-segment (most likely they are for the best known honest block, but don't need to be).
Each $L$ block, in and of itself, only adds $\sim w$ weight to the $L$ chain.
However, if that \emph{draft block} contains $\sim n$ pending PoRs, and each PoR provides $\sim w$ weight on average, then a valid block (created from that draft) will contribute, via PoRs, $\sim nw$ weight to the \emph{honest chain-segment}.
Thus, the miner has this choice: reorganize to the attacker's chain-segment and mine on top of a chain with weight $W_A$, \emph{or} continue mining on the honest chain-segment that weighs $W_H + nw$.

The $nw$ term is that "memory", and it makes all the difference.
Only when $W_A > W_H + nw$ is it safe for an attacker to publish their chain-segment -- then it is \emph{decisively better} than the honest chain-segment.
If, however, the attacker publishes their chain-segment when $W_H < W_A < W_H + nw$: the best chain-segment for honest miners to mine is \emph{still} the honest chain-segment, not the attacker's!

Miners don't actually have to change their behavior at all: their choice is the same, whether they are mining a simplex-chain \emph{or} a traditional chain.
That choice is: \emph{of all possible blocks to mine, which results in the greatest reward?}
If they build on the honest chain-segment, does their resulting block have more total chain-weight than if they were to build on the attacker's chain-segment?

There are some subtleties to consider if chains are DAGs or use GHOST: both chain-segments can be included as ancestors, but which should take priority?
If the weight of PoRs is not taken into account (or applied after evaluating the order of parents), then the attacker's chain-segment takes priority when $W_A > W_H$.
This rule conflicts with what we discussed above, so it must be that parents are ordered based on their chain-weights \emph{including} any PoRs in that block -- the honest chain-segment should take priority when $W_A < W_H + nw$.
Also, note that this is why the weight of PoRs is attributed to the \emph{reflected} block (usually a parent) rather than the block that contains the PoRs.
In essence, the inclusion of this "memory" requires that the fork rule take \emph{context} into account -- particularly the context of the descendant block.
This is not a paradox because the fork rule is \emph{only} used to choose the priority chain in the context of some descendant block (even if that block is an invalid draft).

\subsubsubsection{Testing the Confirmation Equivalence Conjecture}

\label{sec:por-simulation-results}

Let's test the Confirmation Equivalence Conjecture.

\subsubsubsubsection{Hypothesis}

The hypothesis is that: doublespends against an \emph{individual simplex-chain} in an $N$-chain simplex after $C$ \emph{local} confirmations are \emph{as hard} as doublespends against a traditional blockchain after $N\cdot C$ confirmations.

Define $P(N_1 = N; c = C; q = Q)$ (for some given $N$, $C$, and $Q$) to be the probability of an attack succeeding after $c$ confirmations against a simplex with $N_1$ simplex-chains, given an attacker with $q$ proportion of the network hash-rate.

Note that a simplex of 1 chain (the 0-simplex) \emph{is} a traditional blockchain.
So we can take Bitcoin for example: the probability of an attack succeeding after 6 confirmations is given by $P(N_1 = 1; c = 6; q = Q)$ -- of course, we still need to know the attacker's hash-rate to calculate this.

Thus, the hypothesis is that:
\begin{equation}
    P(N_1 = N; c = C; q = Q) \approx P(N_1 = 1; c = NC; q = Q)
    \label{eq:por-sim-hypothesis}
\end{equation}

\subsubsubsubsection{Method}

I have implemented a model blockchain network with attack simulations wherein each chain uses PoW and has a fully functioning fork rule.
This design is modular, such that it supports different fork rules, different block data structures and ordering algorithms (including multiple parents), different attack strategies.
The program also accepts a variety of parameters -- both for the simulation itself, the individual chains, and for each attack strategy.

The source code for this model blockchain network is contained in the \href{https://github.com/AmarooHQ/whitepaper/tree/master/experiments/por-sim-rs}{\texttt{/experiments/por-sim-rs/}} folder of \href{https://github.com/AmarooHQ/whitepaper}{this paper's repository}.

As we can optimize this model, we can simulate many attacks quickly; so we can efficiently make statistically significant measurements.

\todo{asdf - continue method}

\subsubsubsubsection{Scope and Constraints}

This experiment is not intended to be comprehensive -- many significant real-world factors are basically ignored: e.g., block propagation latency, transaction processing overhead, geographic latency, bandwidth constraints, etc.

The scope is \emph{only} intended to cover PoR as a consensus mechanism -- particularly with respect to doublespends.

This constraint means there are many ways to simplify implementation -- for example, reflected blocks are not checked against the history of a blockchain.
These simplifications would render a \emph{production} blockchain insecure.
However, the simplifications are okay in this case since we \emph{expect} near-identical behavior from well functioning networks, and the simplifications are otherwise independent of what we are testing -- the consensus mechanism.
In essence, we want to simulate a \emph{specific} attack, so the implementation doesn't need to be hardened because we know that the attacker will \emph{only} attack in the specific way we care about.

So, while it is \emph{possible} that a poor implementation of a simplex could compromise its security, we'll assume these kinds of design flaws are \emph{in principle} avoidable.

\todoDraftOnly{Test chains of different hash-rates + anything else to cover bases re CEC}

\subsubsubsubsection{Interpreting the Results}

\todo{por sim - Interpreting the Results}

\subsubsubsubsection{Results}

\todo{por sim results}

\subsubsubsubsection{Discussion}

\todo{por sim discussion}

\subsubsubsubsection{Experimental Limitations}

\todo{por sim - limits}

* could be insecure for reasons we didn't test for
* those ignored factors could be decisive
* insufficient breadth of attack parameters
* didn't come up with a good strategy (mb there's a better one than mining in private)
* confirmation bias in implementation -- might have missed something
* does it work for simplexes that have PoS/PoA chains? not tested for
